
\documentclass[12pt,a4paper]{article}

\usepackage[a4paper,margin=1in]{geometry}
\usepackage{graphicx}
\usepackage{booktabs}
\usepackage{amsmath,amssymb}
\usepackage{float}
\usepackage{hyperref}
\usepackage{xcolor}
\usepackage{listings}
\usepackage{enumitem}
\usepackage{fancyhdr}
\usepackage{longtable}

\hypersetup{colorlinks=true,linkcolor=blue,urlcolor=blue}

\pagestyle{fancy}
\fancyhf{}
\lhead{ICS1512}
\rhead{Machine Learning Algorithms Laboratory}
\cfoot{\thepage}

\lstset{
language=Python,
basicstyle=\ttfamily\small,
numbers=left,
frame=single,
breaklines=true,
showstringspaces=false
}

\begin{document}

\begin{tabular}{|p{4cm}|p{11cm}|}
\hline
Experiment No. & 1 \\ \hline
Experiment Title & Working with Python Packages -- NumPy, Pandas, SciPy, Scikit-Learn, Matplotlib
\footnote{\textbf{GitHub Repository: }\url{https://github.com/Barath-j593/ML-assignments}}\\ \hline
Student Name & NavinKumar S \\ \hline
Register Number & 3122247001039 \\ \hline
Faculty & Dr. Poreddy Ajay Kumar Reddy \\ \hline
Submission Date & 16/08/26 \\ \hline
\end{tabular}

\section{Objective}
To explore the core Python libraries used in a Machine Learning workflow -- NumPy, Pandas,
SciPy, Scikit-learn, and Matplotlib -- and to understand key operations such as array
manipulation, data preprocessing, mathematical computing, model building, and data
visualization. Further, to explore publicly available datasets from the UCI Machine Learning
Repository and Kaggle, identify the type of Machine Learning task associated with each, and
perform Exploratory Data Analysis (EDA) on them.

\section{Problem Statement}
Given five publicly available datasets -- Iris, Loan Amount Prediction, Diabetes Prediction,
Email Spam Classification, and MNIST Handwritten Digit Classification -- the task is to load
each dataset, identify the Machine Learning task it corresponds to (regression, binary
classification, or multi-class classification), and perform exploratory data analysis to
understand its structure, quality, and distribution. The input to the workflow is a raw
dataset file (CSV); the output is a statistical and visual understanding of the dataset
(distributions, missing values, correlations, class balance, outliers) that guides the choice
of feature selection technique and ML algorithm for each dataset.

\section{Dataset Description}

\subsection*{Iris Dataset}
\begin{longtable}{|p{6cm}|p{8cm}|}
\hline
Dataset Name & Iris \\ \hline
Dataset Source & UCI Machine Learning Repository (via \texttt{sklearn.datasets.load\_iris}) \\ \hline
Number of Samples & 150 \\ \hline
Number of Features & 4 (sepal length, sepal width, petal length, petal width) \\ \hline
Number of Classes & 3 (setosa, versicolor, virginica -- 50 each) \\ \hline
Missing Values & 0 \\ \hline
Train-Test Split & Not performed -- exploratory analysis only \\ \hline
\end{longtable}

\subsection*{Pima Indians Diabetes Dataset}
\begin{longtable}{|p{6cm}|p{8cm}|}
\hline
Dataset Name & Pima Indians Diabetes \\ \hline
Dataset Source & UCI Machine Learning Repository / Kaggle \\ \hline
Number of Samples & 768 \\ \hline
Number of Features & 8 (Pregnancies, Glucose, BloodPressure, SkinThickness, Insulin, BMI, DiabetesPedigreeFunction, Age) \\ \hline
Number of Classes & 2 (Outcome: 500 non-diabetic, 268 diabetic) \\ \hline
Missing Values & 0 null cells; however 5 Glucose and 374 Insulin values are recorded as 0, which is physiologically implausible and indicates disguised missing data \\ \hline
Train-Test Split & Not performed -- exploratory analysis only \\ \hline
\end{longtable}

\subsection*{Loan Prediction Dataset}
\begin{longtable}{|p{6cm}|p{8cm}|}
\hline
Dataset Name & Loan Prediction \\ \hline
Dataset Source & Kaggle / Analytics Vidhya \\ \hline
Number of Samples & 614 \\ \hline
Number of Features & 11 (excluding Loan\_ID and target Loan\_Status) \\ \hline
Number of Classes & 2 (Loan\_Status: 422 Approved, 192 Rejected) \\ \hline
Missing Values & 149 missing cells across 7 columns (worst: Credit\_History -- 50) \\ \hline
Train-Test Split & Not performed -- exploratory analysis only \\ \hline
\end{longtable}

\subsection*{SMS Spam Collection Dataset}
\begin{longtable}{|p{6cm}|p{8cm}|}
\hline
Dataset Name & SMS Spam Collection \\ \hline
Dataset Source & UCI Machine Learning Repository \\ \hline
Number of Samples & 5,572 \\ \hline
Number of Features & 1 text feature (message content) \\ \hline
Number of Classes & 2 (4,825 ham, 747 spam) \\ \hline
Missing Values & 0 (403 duplicate rows present) \\ \hline
Train-Test Split & Not performed -- exploratory analysis only \\ \hline
\end{longtable}

\subsection*{MNIST Dataset (Sample)}
\begin{longtable}{|p{6cm}|p{8cm}|}
\hline
Dataset Name & MNIST Handwritten Digits (3,000-row sample) \\ \hline
Dataset Source & MNIST database (LeCun et al.), CSV mirror \\ \hline
Number of Samples & 3,000 \\ \hline
Number of Features & 784 (28$\times$28 grayscale pixels) \\ \hline
Number of Classes & 10 (digits 0--9, roughly balanced -- 270 to 332 each) \\ \hline
Missing Values & 0 \\ \hline
Train-Test Split & Not performed -- exploratory analysis only \\ \hline
\end{longtable}

\section{Exploratory Data Analysis}
Include all relevant visualizations.
\begin{itemize}
\item Dataset overview
\item Statistical summary
\item Missing value analysis
\item Class distribution
\item Correlation matrix
\item Feature distribution
\end{itemize}

\subsection*{Iris Dataset}

\begin{figure}[H]
\centering
\includegraphics[width=0.55\linewidth]{iris_species_count.png}
\caption{Species count plot -- Iris}
\end{figure}
\textbf{Inference:} The three species (setosa, versicolor, virginica) are perfectly balanced
at 50 samples each. This balance means accuracy is a reliable evaluation metric here and no
class-imbalance handling (e.g. resampling) is required before training a classifier.

\begin{figure}[H]
\centering
\includegraphics[width=0.6\linewidth]{iris_corr.png}
\caption{Correlation heatmap -- Iris}
\end{figure}
\textbf{Inference:} Petal length and petal width are very strongly correlated with each
other, and both correlate strongly with sepal length, while sepal width correlates weakly
and even negatively with the other features. This indicates that petal measurements alone
carry most of the discriminative signal for species classification, and the two petal
features are partially redundant with one another.

\begin{figure}[H]
\centering
\includegraphics[width=0.85\linewidth]{iris_pairplot.png}
\caption{Pair plot coloured by species -- Iris}
\end{figure}
\textbf{Inference:} Setosa is linearly separable from the other two species on almost every
feature pair, especially petal length vs. petal width, while versicolor and virginica show
some overlap. This suggests a simple linear or distance-based classifier (Logistic
Regression, KNN) should already perform well, with most residual error expected between
versicolor and virginica.

\begin{figure}[H]
\centering
\includegraphics[width=0.45\linewidth]{iris_box_sepal_width.png}
\caption{Sepal width box plot -- Iris}
\end{figure}
\textbf{Inference:} The IQR method flags 4 outlier values in sepal width out of 150 samples
(about 2.7\%). Given the small number and that Iris is a well-curated benchmark dataset,
these are most plausibly genuine measurement variation rather than data-entry errors, so
retaining them is reasonable.

\subsection*{Diabetes Dataset}

\begin{figure}[H]
\centering
\includegraphics[width=0.55\linewidth]{diabetes_hist_glucose.png}
\caption{Glucose distribution -- Diabetes}
\end{figure}
\textbf{Inference:} Glucose is close to a bell-shaped distribution but with a noticeable
spike of zero values -- 5 records have Glucose $=0$, which is physiologically impossible and
indicates the value is actually a missing-data placeholder rather than a true zero reading.

\begin{figure}[H]
\centering
\includegraphics[width=0.55\linewidth]{diabetes_class_count.png}
\caption{Outcome class distribution -- Diabetes}
\end{figure}
\textbf{Inference:} Outcome is imbalanced (500 non-diabetic vs. 268 diabetic, roughly 65:35).
This is mild but non-trivial imbalance, so precision, recall, and F1-score for the positive
class should be reported alongside accuracy, and class weighting or resampling could improve
recall on the minority (diabetic) class.

\begin{figure}[H]
\centering
\includegraphics[width=0.45\linewidth]{diabetes_box_insulin.png}
\caption{Insulin box plot -- Diabetes}
\end{figure}
\textbf{Inference:} The IQR test flags 34 outliers in Insulin. Combined with the fact that
374 of 768 records (about 49\%) have Insulin exactly equal to 0, this again points to
placeholder zeros for unrecorded measurements rather than genuine physiological readings --
this column will need explicit missing-value imputation rather than outlier removal.

\begin{figure}[H]
\centering
\includegraphics[width=0.75\linewidth]{diabetes_corr.png}
\caption{Correlation heatmap -- Diabetes}
\end{figure}
\textbf{Inference:} Glucose shows the strongest correlation with Outcome (Pearson $r \approx
0.47$) of any single feature, consistent with its clinical role in diabetes diagnosis. Other
features (BMI, Age, Pregnancies) show weaker but still positive associations, while no pair
of predictors is correlated highly enough to suggest serious multicollinearity.

\subsection*{Loan Prediction Dataset}

\begin{figure}[H]
\centering
\includegraphics[width=0.7\linewidth]{loan_missing.png}
\caption{Missing value heatmap -- Loan Prediction}
\end{figure}
\textbf{Inference:} Missing values are spread across seven columns, most severely
Credit\_History (50 missing), Self\_Employed (32), and LoanAmount (22). Credit\_History is
typically the single strongest predictor of loan approval, so its missingness needs careful
handling (mode/most-frequent imputation, or a separate ``unknown'' category) rather than
row deletion, which would discard a meaningful fraction of the 614 records.

\begin{figure}[H]
\centering
\includegraphics[width=0.55\linewidth]{loan_hist_income.png}
\caption{Applicant income distribution -- Loan Prediction}
\end{figure}
\textbf{Inference:} Applicant income is heavily right-skewed (skewness $\approx 6.5$), with
most applicants clustered at lower incomes and a long tail of high earners. This strongly
suggests a log transformation before feeding income into a regression model, and explains
why 50 values are flagged as IQR outliers on the raw scale.

\begin{figure}[H]
\centering
\includegraphics[width=0.55\linewidth]{loan_status_count.png}
\caption{Loan status distribution -- Loan Prediction}
\end{figure}
\textbf{Inference:} Loan\_Status is imbalanced, with about 69\% of applications approved (Y:
422) versus 31\% rejected (N: 192). This imbalance should be reflected in the evaluation
metric choice (F1/recall on the rejected class matter as much as overall accuracy) if the
downstream task is framed as approval classification rather than amount regression.

\begin{figure}[H]
\centering
\includegraphics[width=0.6\linewidth]{loan_property_area.png}
\caption{Property area distribution -- Loan Prediction}
\end{figure}
\textbf{Inference:} Applications are reasonably distributed across the three property-area
categories (Urban, Semiurban, Rural), with no single category dominating -- so one-hot
encoding this feature should not introduce severe sparsity or imbalance concerns.

\subsection*{Email Spam Dataset}

\begin{figure}[H]
\centering
\includegraphics[width=0.5\linewidth]{spam_class_count.png}
\caption{Ham vs. spam class distribution}
\end{figure}
\textbf{Inference:} The dataset is strongly imbalanced -- 4,825 ham messages versus only 747
spam (about 87:13). A classifier trained naively on raw accuracy could reach $\sim$87\% by
always predicting ``ham,'' so precision/recall/F1 on the spam class, or techniques such as
class weighting, are essential for a meaningful evaluation.

\begin{figure}[H]
\centering
\includegraphics[width=0.55\linewidth]{spam_char_len.png}
\caption{Character length distribution -- Spam}
\end{figure}

\begin{figure}[H]
\centering
\includegraphics[width=0.55\linewidth]{spam_word_count.png}
\caption{Word count distribution -- Spam}
\end{figure}
\textbf{Inference:} Message length is itself discriminative: spam messages average about 139
characters versus roughly 71 for ham -- almost double. Both distributions are right-skewed
with a long tail toward longer messages, so character/word count could be added as an
engineered feature alongside TF-IDF text features.

\subsection*{MNIST Dataset}

\begin{figure}[H]
\centering
\includegraphics[width=0.55\linewidth]{mnist_class_count.png}
\caption{Digit class distribution -- MNIST}
\end{figure}
\textbf{Inference:} The 10 digit classes are close to balanced in this sample, ranging from
270 (digit 5) to 332 (digit 1) examples -- roughly a 20\% spread. This near-balance means
accuracy remains a reasonable primary metric, though a per-class confusion matrix is still
useful given digits like 4/9 or 3/8 are visually similar and prone to confusion.

\begin{figure}[H]
\centering
\includegraphics[width=0.7\linewidth]{mnist_samples.png}
\caption{Sample digit images -- MNIST}
\end{figure}
\textbf{Inference:} The sample grid confirms the images are correctly labelled, centered,
grayscale $28\times28$ handwritten digits consistent with the standard MNIST format, with
visible natural variation in stroke width and slant across samples of the same digit.

\begin{figure}[H]
\centering
\includegraphics[width=0.55\linewidth]{mnist_pixel_dist.png}
\caption{Pixel intensity distribution -- MNIST}
\end{figure}
\textbf{Inference:} Pixel values range from 0 to 255 with a mean of about 33 and a large
standard deviation ($\approx$78), and only about 19\% of pixels are non-zero. This confirms
the images are mostly dark background with sparse bright strokes, and supports normalizing
pixel values (dividing by 255) before feeding them into a CNN or any distance-based model.

\section{Data Preprocessing}

Explain:
\begin{itemize}
\item \textbf{Missing value handling:} Iris and MNIST require no imputation. Diabetes needs
the disguised zero-placeholders in Glucose and Insulin replaced with median/mean imputation
per Outcome class. Loan needs mode imputation for categorical columns (Gender,
Self\_Employed, Credit\_History) and median imputation for LoanAmount.
\item \textbf{Encoding:} Categorical features in Loan (Gender, Married, Education,
Self\_Employed, Property\_Area) and the target labels in Diabetes/Spam/Iris/Loan require
Label or One-Hot Encoding before modelling.
\item \textbf{Feature scaling:} Diabetes and Loan numeric features (Glucose, BMI,
ApplicantIncome) span very different ranges and should be standardized
(\texttt{StandardScaler}) or normalized (\texttt{MinMaxScaler}). MNIST pixel values should be
scaled to $[0,1]$ by dividing by 255.
\item \textbf{Train-test split:} Not performed in this experiment, which is limited to
exploratory analysis; a typical 80:20 or 70:30 stratified split would be used in a subsequent
modelling experiment.
\item \textbf{Feature engineering (if applicable):} For Spam, message character length and
word count (already explored in Section 4) could be added as engineered numeric features
alongside text-vectorized (TF-IDF) features.
\end{itemize}

\section{Machine Learning Algorithm}

Based on the EDA findings, the following algorithms are recommended for each dataset (feature
selection technique and full task identification is summarized in the table below).

\begin{longtable}{|p{3.3cm}|p{2.6cm}|p{3.6cm}|p{4.7cm}|}
\hline
\textbf{Dataset} & \textbf{Type of ML Task} & \textbf{Feature Selection Technique} & \textbf{Suitable ML Algorithm} \\ \hline
\endhead
Iris & Supervised -- Multi-class Classification & ANOVA F-test (\texttt{SelectKBest} with \texttt{f\_classif}) & Logistic Regression / K-Nearest Neighbors / Decision Tree \\ \hline
Loan Amount Prediction & Supervised -- Regression & Correlation analysis, ANOVA F-test (\texttt{f\_regression}) & Linear Regression / Random Forest Regressor \\ \hline
Predicting Diabetes & Supervised -- Binary Classification & Chi-square test (\texttt{SelectKBest} with \texttt{chi2}) & Logistic Regression / Random Forest Classifier \\ \hline
Classification of Email Spam & Supervised -- Binary Classification (Text) & Chi-square test on TF-IDF/Bag-of-Words features & Naive Bayes / Support Vector Machine (SVM) \\ \hline
Handwritten Character Recognition / MNIST & Supervised -- Multi-class Classification (Image) & Variance thresholding / PCA (dimensionality reduction over pixel features) & Convolutional Neural Network (CNN) / SVM / Random Forest \\ \hline
\end{longtable}

\textbf{Working principle (representative example -- Logistic Regression):} Models the
probability of a binary/multi-class outcome using a linear combination of input features
passed through a sigmoid (or softmax) function, trained by minimizing log-loss via gradient
descent.

\textbf{Mathematical formulation:}
$$P(y=1\mid x) = \sigma(w^\top x + b) = \frac{1}{1+e^{-(w^\top x + b)}}$$

\textbf{Advantages:} Simple, interpretable, computationally cheap, works well when classes
are near-linearly separable (as seen for setosa in the Iris pair plot).

\textbf{Limitations:} Struggles with non-linear decision boundaries and is sensitive to
unscaled features and multicollinearity -- both relevant given the skewed income feature in
the Loan dataset and the correlated petal measurements in Iris.

\section{Experimental Setup}
\begin{tabular}{ll}
\toprule
Python Version & 3.11.9\\
Scikit-Learn Version & 1.8.0\\
IDE & Jupyter Notebook \\
Operating System & Window\\
Hardware &  Victus by HP Gaming Laptop 15-fa2xxx\\
\bottomrule
\end{tabular}

\section{Hyperparameter Selection}
Not applicable for this experiment -- Experiment 1 covers library exploration, dataset
identification, and EDA only; no model was trained or tuned.

\begin{tabular}{ll}
\toprule
Parameter & Value\\
\midrule
-- & -- \\
\bottomrule
\end{tabular}

\section{Results}
Not applicable for this experiment -- no model was trained. The equivalent deliverable for
this experiment, the ML task/feature-selection/algorithm identification table, is provided in
Section 6.

\begin{tabular}{lc}
\toprule
Metric & Value\\
\midrule
Accuracy & -- \\
Precision & -- \\
Recall & -- \\
F1-score & -- \\
ROC-AUC & -- \\
\bottomrule
\end{tabular}

\section{Result Visualization}
Not applicable for this experiment -- no model was trained, so no confusion matrix, ROC
curve, or learning curve is available. The dataset-level visualizations (distributions,
correlations, class balance) generated for this experiment are provided in Section 4.

\section{Discussion}
Across all five datasets, EDA surfaced issues that would materially affect modelling: disguised
missing values (zeros) in Diabetes, genuine missing values concentrated in a few key columns
in Loan, class imbalance in Diabetes/Loan/Spam, and a heavily right-skewed income feature in
Loan. Iris and MNIST were the cleanest datasets, with no missing values and reasonably
balanced classes. These findings directly justify the feature selection technique and
algorithm choices recommended in Section 6 -- e.g. chi-square/ANOVA tests for the tabular
datasets, TF-IDF-based selection for text, and PCA/variance thresholding for the
high-dimensional MNIST pixel features.

\section{Conclusion}
This experiment explored the core NumPy, Pandas, SciPy, Scikit-learn, and Matplotlib/Seaborn
functionality needed for a Machine Learning workflow, and used it to identify the ML task
type, suitable feature selection technique, and suitable algorithm for five public datasets
(Iris, Loan Amount Prediction, Diabetes Prediction, Email Spam Classification, and MNIST /
Handwritten Character Recognition). Exploratory data analysis on the real datasets revealed
concrete data-quality issues (disguised missing values, skewed distributions, class
imbalance) that will directly inform preprocessing decisions in subsequent experiments.

\section{References}
\begin{enumerate}[label={[\arabic*]}]
\item R. A. Fisher, ``The Use of Multiple Measurements in Taxonomic Problems,'' \textit{Annals of Eugenics}, 1936. Dataset accessed via \url{https://scikit-learn.org/stable/datasets/toy_dataset.html}
\item National Institute of Diabetes and Digestive and Kidney Diseases, ``Pima Indians Diabetes Dataset,'' UCI Machine Learning Repository, \url{https://archive.ics.uci.edu/}
\item Analytics Vidhya, ``Loan Prediction Dataset,'' Kaggle, \url{https://www.kaggle.com/datasets}
\item T. A. Almeida, J. M. G. Hidalgo, ``SMS Spam Collection Dataset,'' UCI Machine Learning Repository, \url{https://archive.ics.uci.edu/}
\item Y. LeCun, C. Cortes, C. J. C. Burges, ``The MNIST Database of Handwritten Digits,'' \url{http://yann.lecun.com/exdb/mnist/}
\item Scikit-learn Documentation, \url{https://scikit-learn.org/stable/}
\item Pandas Documentation, \url{https://pandas.pydata.org/docs/}
\end{enumerate}

\end{document}

